\section*{A060: exact selector-tax sharpness for a Boolean point defect}

All Fourier coefficients are normalized on the uniform cube and all
logarithms are base two.  Define the genuine Boolean functions
\[
 D(z_1,z_2)=1-2\mathbf 1_{\{z_1=z_2=-1\}},\qquad
 K(x,z_1,z_2)=1-2\mathbf 1_{\{x=z_1=z_2=-1\}}.
\]
Equivalently, with the fresh selector convention (x=+1) for the first
branch,
\[
 K(+1,z)=1,\qquad K(-1,z)=D(z).
\]

\paragraph{Proposition (PROVED, exact finite witness).}
For the splice above,
\[
 \frac{H(K)-\tfrac12 H(D)}{I_x(K)}
 =12-\frac94\log_2 9
 \in(4.867,4.868).
\]
Consequently, every proposed inequality valid for all constant-versus-Boolean
one-bit splices of the form
\[
 H(K)\leq \frac12 H(D)+c I_x(K)
\]
must have (c\geq 12-\frac94\log_2 9).  In particular (c=4) is false,
already for a genuine three-bit Boolean function.  This is a lower bound on
this local scalar selector tax, not a lower bound of the same size on the
universal FEI constant.

\paragraph{Proof.}
Direct expansion gives
\[
 D=\frac12+\frac12z_1+\frac12z_2-\frac12z_1z_2.
\]
Thus its four squared Fourier masses are all (1/4), whence
\[
 H(D)=2,\qquad I(D)=1,\qquad \mathbb E D=\frac12.
\]
The three-bit point defect (K) has constant Fourier coefficient (3/4)
and each of its seven nonconstant Fourier coefficients has magnitude (1/4).
Therefore
\[
 H(K)=\frac9{16}\log_2\frac{16}{9}+\frac74,
 \qquad I(K)=\frac34.
\]
Exactly four of the seven nonconstant characters contain (x), so
\[
 I_x(K)=4\cdot\frac1{16}=\frac14.
\]
It follows that
\begin{align*}
 \frac{H(K)-\tfrac12H(D)}{I_x(K)}
 &=4\left(\frac9{16}\log_2\frac{16}{9}+\frac34\right)\\
 &=12-\frac94\log_2 9.
\end{align*}
The attached rational-interval certificate proves that this number lies in
((4.867,4.868)), without floating-point assumptions.  This proves the
claim. \qed

\paragraph{Actual FEI ratio and tensor scope.}
The same exact ledger yields
\[
 \frac{H(K)}{I(K)}
 =\frac{16}{3}-\frac34\log_2 9
 \in(2.955,2.956).
\]
For the disjoint tensor product (K^{\otimes m}), the squared Fourier law is
the (m)-fold product law.  Hence Shannon entropy and total influence are
both additive:
\[
 H(K^{\otimes m})=mH(K),\qquad
 I(K^{\otimes m})=mI(K).
\]
The tensor family therefore preserves the smaller ratio
(16/3-(3/4)\log_2 9); it does not promote the local tax coefficient
(12-(9/4)\log_2 9) into an FEI lower bound.

\paragraph{Correlation and hypothesis audit.}
The two branches are the constant (1) and (D), so their Fourier inner
product is (\langle 1,D\rangle=\mathbb E D=1/2).  The example uses a fresh
selector but is not neighboring-orthogonal.  It has no reused selector and no
shared data inputs between separate children.  It is therefore a sharpness
test for the nonorthogonal scalar overhead in a fresh-selector splice, not an
overlapping-input theorem.  All gates are Boolean; there is no weighted or
non-Boolean relaxation.

\paragraph{Endpoint and guard checks.}
The witness is not a parity-zero erasure: both branch means and the branch
correlation are explicit.  There is no rare-event asymptotic or hidden
dimension limit.  The disjoint tensor statement uses only exact product
factorization and does not assert a shared-variable product law.  The result
neither contradicts disjoint additivity nor repairs the registered
\texttt{OVERLAP\_PRODUCT} obstruction.

